\begin{lemma}[Presentation cases]
The structural equivalence of \noderef{PowerPresentationEq} has the three
exclusive shapes required by the hierarchy.  No atom is equivalent to a
node, in either order.  For nodes
$\mathsf{node}(f)$ and $\mathsf{node}(g)$, equivalence is exactly the
conjunction
\[
 (\forall i\,\exists j\; f(i)\mathrel{\mathsf{PowerPresentationEq}}g(j))
 \ \land\ 
 (\forall j\,\exists i\; f(i)\mathrel{\mathsf{PowerPresentationEq}}g(j)).
\]
\end{lemma}
\begin{proof}
Unfold the four defining equations in
\noderef{PowerPresentationEq}.  The atom--node and node--atom equations
reduce to $\mathsf{False}$, proving both disjointness assertions.  The
node--node equation is precisely the displayed forward and backward
matching conjunction, while the atom--atom equation is equality of the
underlying $Q$-values.  These cases exhaust the two tagged constructors of
\noderef{PowerQ}, so the stated characterization is complete.
\end{proof}
